% paper/sections/02_related_work.tex
\section{Related Work}
\label{sec:related}

\subsection{Supervised LLM Routing}

\citet{ong2024routellm} propose RouteLLM, which trains a classifier on 55K Chatbot Arena human preference labels. Their Similarity-Weighted (SW) router retrieves nearest neighbors from the arena embedding space and uses Elo-weighted win rates as routing scores. While effective when pretraining data aligns with deployment distribution, SW requires retraining when model capabilities change and is architecturally binary. Not Diamond \citep{notdiamond2024} uses a random forest (RoRF) over benchmark performance features, also requiring pretraining labels per model pair.

\citet{chen2023frugalgpt} propose FrugalGPT, a cascade system that queries a sequence of models in increasing cost order, stopping when the response quality exceeds a threshold. \citet{dekoninck2025unified} provide a unified framework for routing and cascading, proving theoretical guarantees for cascade-based approaches. These methods avoid per-query pretraining but require a quality estimator and a cascade threshold.

\citet{shnitzer2023large} study mixture-of-experts routing for LLMs but focus on fine-tuned specialist models rather than black-box API routing. \citet{lu2023routing} use embedding similarity to route queries to domain-specialist models.

\subsection{Online Learning for LLM Serving}

LinUCB \citep{li2010contextual} is the canonical contextual bandit algorithm for recommendation and routing. \citet{russo2018tutorial} provide a tutorial on Thompson Sampling, showing it matches LinUCB's regret while requiring less tuning. \citet{abeille2017linear} prove $O(d\sqrt{T\log T})$ regret for Linear Thompson Sampling.

\subsection{Multi-Model LLM Systems}

\citet{aggarwal2024automix} propose AutoMix, which uses a POMDP to decide when to escalate from a weak to a strong model. LLM-Blender \citep{jiang2023llmblender} ensembles outputs from multiple models using pairwise ranking, but requires all models to respond to each query.

\textbf{Comparison to RouteSmith:}

RouteSmith is a multi-strategy system: its online bandit component (this paper) learns from per-query feedback with zero pretraining, while its random forest predictor \citep{routesmith2025system} provides a supervised alternative for cold-start and calibrated deployment. This dual architecture lets practitioners choose between zero-shot online learning and offline pretraining depending on their dataset availability.

\begin{table}[h]
\centering
\begin{tabular}{lcccc}
\toprule
Method & Online & N-arm & No labels & Interpretable \\
\midrule
RouteLLM-SW & \xmark & \xmark & \xmark & \xmark \\
Not Diamond (RoRF) & \xmark & \xmark & \xmark & \xmark \\
FrugalGPT & \cmark & \cmark & \xmark & \xmark \\
Dekoninck et al.\ (cascade) & \cmark & \cmark & \xmark & \xmark \\
TS-Cat & \cmark & \cmark & \cmark & \xmark \\
LinUCB-27d & \cmark & \cmark & \cmark & \cmark \\
\textbf{LinTS-27d (Ours)} & \cmark & \cmark & \cmark & \cmark \\
\bottomrule
\end{tabular}
\caption{RouteSmith vs. prior work on four dimensions. Not Diamond's RoRF is included as a prominent open-source static router; it shares the same structural limitations as RouteLLM-SW. Dekoninck et al.\ (cascade) represents the state of the art in heuristic cascade routing with theoretical guarantees.}
\end{table}
